% 4.6 — Le côté Windows et C3
\section{Le côté Windows}
\label{sec:windows}

\subsection{Parité entre les deux systèmes}
Le tableau~\ref{tab:parite} récapitule, couche par couche, ce qui est identique et ce qui
diffère.

\begin{table}[htbp]
\centering
\caption{Parité Linux / Windows au commit de référence}
\label{tab:parite}
\small
\begin{tabularx}{\linewidth}{@{}L{3.2cm}YY@{}}
\toprule
\textbf{Couche} & \textbf{Linux} & \textbf{Windows} \\
\midrule
Portes & \textit{hook} Claude Code (Bash, fichiers), Codex, Cursor, piège bash &
  \textit{hook} Claude Code (Bash/Git Bash, PowerShell, fichiers), Codex (PowerShell) \\
Lecture POSIX & moteur & même moteur ; chemins \cmd{/c/x} traduits en
  \cmd{C:\textbackslash x} \\
Lecture PowerShell & --- & C3, même moteur ; les anciennes règles ne gardent que leurs captures et
  escalades \\
Politique, décideur, reçus, \cmd{undo} & communs & communs \\
Filet d'exécution & garde d'appels système (ptrace), activé par réécriture &
  montage WinFsp : le projet \emph{est} le montage, aucune réécriture \\
Plan de contrôle pendant l'exécution & appel annulé par le garde &
  refusé par les portes (le montage lui-même : V7, ouvert) \\
Persistance & --- & tâche d'ouverture de session durcie (W-SEC) \\
\bottomrule
\end{tabularx}
\end{table}

\subsection{Une leçon de la partie T : juger ce qui s'exécute}
Pour faire passer les lignes POSIX de Windows par le moteur, une première version (M2)
\emph{réécrivait le texte} de la commande, en traduisant ses chemins, avant de la juger. La
partie T l'a retirée : bash exécutait alors une autre commande que celle qui avait été jugée,
et les reçus enregistraient un texte que l'agent n'avait jamais écrit. La traduction se fait
désormais à l'intérieur du moteur, sur les chemins, et la commande exécutée reste celle de
l'agent. Le principe est général : \emph{on juge exactement ce qui s'exécute, et on enregistre
exactement ce qui a été jugé}.

\subsection{PowerShell avant C3}
Sous Windows, les agents écrivent du PowerShell. Jusqu'à C3, le moteur ne le lit pas : une ligne
PowerShell, ou une ligne POSIX qui contient un mot PowerShell ou \cmd{cmd.exe}, est jugée par
une poignée d'expressions régulières (\cmd{ps\_remove\_item}, \cmd{ps\_set\_content}\ldots ;
arbitrage M5). Le coût était mesurable (tableau~\ref{tab:ps-avant}) : des destructions passaient,
une écriture banale escaladait, et seul le montage WinFsp couvrait ces cas. C3 les lit
(colonne « Avec C3 »).

\begin{table}[htbp]
\centering
\caption{Lignes PowerShell : l'ancien chemin et C3 (le projet contient a.txt et b.txt)}
\label{tab:ps-avant}
\small
\begin{tabularx}{\linewidth}{@{}L{4.4cm}L{3.9cm}Y@{}}
\toprule
\textbf{Ligne} & \textbf{Ancien chemin} & \textbf{Avec C3} \\
\midrule
\cmd{iex "Remove-Item a.txt"} & ALLOW (le motif voit une chaîne) & REVERSIBLE : la chaîne
  littérale est analysée, a.txt est copié avant \\
\cmd{Write-Output "\$(Remove-Item a.txt)"} & ALLOW (suppression cachée) & REVERSIBLE : la
  sous-expression est analysée \\
\cmd{Copy-Item a.txt b.txt} & ALLOW (écrase sans \cmd{-Force}) & REVERSIBLE : b.txt, la
  destination écrasée, est copié avant \\
\cmd{Set-Content a.txt hi} & ESCALATE (inutile) & REVERSIBLE : a.txt lié à \cmd{-Path}
  (position 0), hi à \cmd{-Value} \\
\cmd{"a.txt" | Remove-Item} & ALLOW (la cible arrive par le tube) & REVERSIBLE : la valeur se
  lie au paramètre de \textit{pipeline} \\
\bottomrule
\end{tabularx}
\end{table}

\section{C3 : le second frontal du moteur}
\label{sec:c3}

\subsection{L'idée : un frontal, pas un second garde-fou}
Une expression régulière lit des caractères ; PowerShell \emph{lie des paramètres}. Seul un
analyseur qui les lie comme PowerShell le fait peut dire ce qu'une ligne touchera. C3 n'est
pas un second garde-fou à côté du premier : c'est un \emph{second frontal} du même moteur. La
ligne PowerShell est analysée vers le même modèle d'invocations, simulée sur la même vue
virtuelle du système de fichiers, et produit les mêmes effets, qui passent par les mêmes
tables \cmd{verdict} et \cmd{policy} et le même décideur (figure~\ref{fig:c3}). Une politique
et un ensemble de garanties --- « jamais un REVERSIBLE non mérité », « l'inconnu n'est pas
sûr » --- valent pour les deux dialectes.

\begin{figure}[htbp]
\centering
% Figure 5b — C3 : deux frontaux, un seul moteur.
\begin{tikzpicture}[
    pas/.style={boite,text width=20mm,font=\sffamily\scriptsize,minimum height=11mm},
    part/.style={accent,text width=16mm,font=\sffamily\scriptsize,minimum height=11mm},
  ]
  \node[douce,text width=15mm,font=\sffamily\scriptsize] (lp) at (0,0.9) {Ligne POSIX};
  \node[douce,text width=15mm,font=\sffamily\scriptsize] (lw) at (0,-0.9) {Ligne PowerShell};
  \node[pas] (pp) at (2.5,0.9) {Analyseur POSIX\sous{\texttt{model.py}}};
  \node[pas] (pw) at (2.5,-0.9) {Analyseur PS\\+ lieur de\\paramètres\sous{C3}};
  \node[part] (inv) at (5.0,0) {Invocations};
  \node[part,right=3mm of inv] (sim) {Simulation\sous{FS virtuel}};
  \node[part,right=3mm of sim] (eff) {Effets\sous{DELETE, MOVE\ldots}};
  \node[part,right=3mm of eff] (pol) {\texttt{verdict} / \texttt{policy}};
  \node[boite,draw=encre,line width=1pt,text width=17mm,font=\sffamily\scriptsize,right=3mm of pol,minimum height=11mm] (dec) {\textbf{\texttt{decide}}\sous{décision + reçu}};
  \draw[fleche] (lp) -- (pp); \draw[fleche] (lw) -- (pw);
  \draw[fleche] (pp.east) -- (inv.west); \draw[fleche] (pw.east) -- (inv.west);
  \foreach \a/\b in {inv/sim,sim/eff,eff/pol,pol/dec} \draw[fleche] (\a) -- (\b);
  \node[zone,draw=wepwnfonce,fit=(inv)(pol),label={[etiquette,text=wepwnfonce]above:partagé par les deux dialectes}] {};
\end{tikzpicture}

\caption{C3 : deux frontaux, un seul moteur}
\label{fig:c3}
\end{figure}

\subsection{Fonctionnement}
\begin{description}[style=nextline,leftmargin=0.6cm,font=\sffamily\bfseries]
  \item[Analyse lexicale et syntaxique]
    Les règles propres à PowerShell : citations simples littérales et doubles qui
    s'étendent, échappement et continuation par l'accent grave, \textit{here-strings}
    \cmd{@'\ldots'@} et \cmd{@"\ldots"@}, \textit{pipelines}, séparateurs \cmd{;},
    sous-expressions \cmd{\$(\ldots)}, blocs de script \cmd{\{\ldots\}}, opérateur d'appel
    \cmd{\&}, redirections, variables \cmd{\$env:} et affectations sur la même ligne.
  \item[Lieur de paramètres]
    Le cœur de C3. Pour chaque cmdlet, il connaît les \emph{jeux de paramètres}, les paramètres
    \emph{positionnels} et leur position, les alias de la cmdlet et de chaque paramètre, les
    paramètres commutateurs, et ceux qui acceptent une entrée de \textit{pipeline} (par valeur
    ou par nom de propriété). Il résout les abréviations comme PowerShell : le plus court
    préfixe non ambigu (\cmd{-Rec} $\rightarrow$ \cmd{-Recurse}, \cmd{-Dest} $\rightarrow$
    \cmd{-Destination}). Ces tables sont \emph{générées} à partir du PowerShell installé
    (\cmd{Get-Command}, métadonnées des paramètres, \cmd{Get-Alias}), et chaque fait de liaison
    a été vérifié sur PowerShell 5.1. Ainsi \cmd{Set-Content a.txt hi} lie \cmd{a.txt} à
    \cmd{-Path} (position 0) et \cmd{hi} à \cmd{-Value} (position 1).
  \item[Effets par cmdlet]
    \cmd{Remove-Item} et ses alias $\rightarrow$ DELETE, récursif avec \cmd{-Recurse} ;
    \cmd{-Path} développe les jokers, \cmd{-LiteralPath} jamais. \cmd{Clear-Content}
    $\rightarrow$ TRUNCATE. \cmd{Set-Content} et \cmd{Out-File} $\rightarrow$ OVERWRITE d'un
    fichier existant, CREATE sinon ; \cmd{-Append} et \cmd{Add-Content} ajoutent seulement ;
    \cmd{-NoClobber} refuse. \cmd{Copy-Item} $\rightarrow$ OVERWRITE de la destination si elle
    existe ; \cmd{Move-Item} et \cmd{Rename-Item} $\rightarrow$ MOVE, et refusent une
    destination existante sans \cmd{-Force}. \cmd{New-Item} $\rightarrow$ CREATE, ou OVERWRITE
    avec \cmd{-Force}. Les lectures (\cmd{Get-ChildItem}, \cmd{Get-Content},
    \cmd{Select-Object}\ldots) n'ont pas d'effet. Avec \cmd{-WhatIf}, PowerShell ne fait que
    décrire l'action : aucun effet.
  \item[\textit{Pipelines}]
    Une valeur littérale envoyée dans une cmdlet se lie à son paramètre de \textit{pipeline}
    (\cmd{"a.txt" | Remove-Item} supprime \cmd{a.txt}) ; \cmd{Get-ChildItem *.tmp |
    Remove-Item} supprime ce que la liste renvoie, et la liste est simulée sur le système de
    fichiers virtuel.
  \item[Du code caché dans du code]
    Une sous-expression est analysée et ses effets comptent ; un bloc de script est analysé là
    où il s'exécute (\cmd{ForEach-Object \{\ldots\}}). \cmd{Invoke-Expression} (\cmd{iex}) avec
    une chaîne \emph{littérale} est analysé comme une ligne à part entière --- l'équivalent de
    \cmd{eval} ou \cmd{bash -c} --- et \cmd{powershell.exe -EncodedCommand} est décodé (UTF-16LE)
    puis analysé. Une chaîne construite à l'exécution ne peut pas être lue : elle est non
    comprise.
  \item[Rien n'est deviné]
    Une construction hors de la grammaire modélisée, une variable de valeur inconnue, une
    cmdlet ou un paramètre inconnus rendent la ligne non comprise ou non couverte, exactement
    comme en POSIX : filet WinFsp, sinon point de contrôle, sinon escalade. Un script lancé par
    \cmd{-File script.ps1} ou \cmd{\& .\textbackslash x.ps1} est un exécuteur de code, traité
    comme \cmd{python x.py}.
\end{description}

Une ligne \cmd{cmd /c \ldots} n'est pas du PowerShell : elle garde les règles textuelles
(arbitrage PS3). Les anciennes règles PowerShell restent en place, mais ne gardent que leurs
captures et leurs escalades, jamais leur ALLOW : la réponse la plus stricte l'emporte. C3 ne lit
pas les chaînes construites à l'exécution, les classes, les fonctions, les modules ni les
scripts \cmd{.ps1}, qui sont des exécuteurs de code.

\subsection{Les arbitrages PS1--PS5}
\begin{itemize}
  \item \textbf{PS1} --- un analyseur PowerShell écrit en Python, dans le processus du
    \textit{hook} ; l'analyseur de PowerShell lui-même sert d'oracle dans les tests sous Windows.
  \item \textbf{PS2} --- les tables du lieur de paramètres sont générées depuis le PowerShell
    installé (5.1.19041 : \cthreecmdlets{} cmdlets, \cthreealias{} alias), pas écrites à la main.
  \item \textbf{PS3} --- \cmd{cmd /c \ldots} garde les règles textuelles.
  \item \textbf{PS4} --- \cmd{-WhatIf} n'a aucun effet : PowerShell décrit l'action sans la
    faire.
  \item \textbf{PS5} --- une livraison en étapes : (a) analyseur, lieur, cmdlets de fichiers et
    lectures, branchement au moteur ; (b) \textit{pipelines}, \cmd{\$( )}, blocs de script,
    \cmd{iex}, \cmd{-EncodedCommand}, \cmd{\&}, \cmd{-File} ; (c) retrait des règles textuelles
    PowerShell, en perspective.
\end{itemize}

\subsection{Validation}
Chaque règle de cmdlet porte ses exemples, exécutés comme tests ; les lignes PowerShell réelles
trouvées dans les registres sont figées avec leur décision attendue ; chaque trou de V12 a un
test qui échoue sur le code d'avant C3 et passe après. Sous Windows, l'analyse de C3 est
comparée, sur un corpus de lignes réelles, à l'arbre syntaxique produit par l'analyseur de
PowerShell lui-même. Le portail de \nbcorpus{} commandes et le portail Windows ont été rejoués :
chaque ligne qui bouge est plus stricte ou arbitrée. Les anciennes règles PowerShell restent en
place jusqu'à ce que les décisions de C3 les égalent ou les dépassent sur chaque ligne du
portail.
